%% LyX 2.3.6.1 created this file.  For more info, see http://www.lyx.org/.
%% Do not edit unless you really know what you are doing.
\documentclass[12pt,a4paper,oneside,english]{amsart}
\usepackage{ae,aecompl}
\renewcommand{\familydefault}{\sfdefault}
\usepackage[T1]{fontenc}
\usepackage[latin9]{luainputenc}
\usepackage{calc}
\usepackage{amsthm}
\usepackage{setspace}
\usepackage[authoryear]{natbib}
\onehalfspacing

\makeatletter

%%%%%%%%%%%%%%%%%%%%%%%%%%%%%% LyX specific LaTeX commands.
% Backwards compatibility for LuaTeX < 0.90
\@ifundefined{pageheight}{\let\pageheight\pdfpageheight}{}
\@ifundefined{pagewidth}{\let\pagewidth\pdfpagewidth}{}
\pageheight\paperheight
\pagewidth\paperwidth


%%%%%%%%%%%%%%%%%%%%%%%%%%%%%% Textclass specific LaTeX commands.
\numberwithin{equation}{section}
\numberwithin{figure}{section}
\theoremstyle{plain}
\newtheorem*{thm*}{\protect\theoremname}
\theoremstyle{definition}
\newtheorem*{defn*}{\protect\definitionname}
\theoremstyle{plain}
\newtheorem{thm}{\protect\theoremname}

\makeatother

\usepackage{babel}
\providecommand{\definitionname}{Definition}
\providecommand{\theoremname}{Theorem}

\begin{document}
\title{Notes on An Introduction to Modern Cosmology by Andrew Liddle}
\maketitle

\section{Overview}

Everything comes from this book.\footnote{Liddle, Andrew. \textit{An Introduction to Modern Cosmology}, 3rd
Edition. John Wiley \& Sons, 2015. \textcompwordmark{}}

\subsection{Expansion of the Universe}

Most of my current knowledge of physics comes from the two courses
I took during my U1 fall semester: PHYS 260 - Modern Physics, and
PHYS 251 Honours Classical Mechanics. I will most likely tend to relate
the material seen in this book to the material I have seen in class. 
\begin{quote}
Firstly, the relativisitc doppler effect formula for a receding source:
\end{quote}
\[
\begin{aligned}f_{obs}= & f_{em}\sqrt{\frac{1-\frac{v}{c}}{1+\frac{v}{c}}}\end{aligned}
\]

Which means that the frequency perceive by the observer from a receding
source is lower, which means that the wavelength is longer, which,
consequently, means that the light is red shifted. The standard terminology
for redshift is:

\[
\begin{aligned}z= & \frac{\lambda_{obs}-\lambda_{em}}{\lambda_{em}}\end{aligned}
\]

where $\lambda_{em}$and $\lambda_{obs}$ are the wavelengths of light
at the points of emission. If a nearby object is receding at speed
v, then its redshift is:

\[
\begin{aligned}z= & \frac{v}{c}\end{aligned}
\]

\begin{thm*}
The velocity of recession is proportional to the distance of the object. 

\[
\vec{v}=H_{0}\overrightarrow{r}
\]

This is known as Hubble's law, and the constant of proportionality
$H_{0}$, which is known as the Hubble's constant.
\end{thm*}
%
This does not violate the Cosmological Principle since \textit{every
}observer sees all object receing away from them with velocity proportional
to distance. 
\begin{defn*}
The Cosmological Principle: The Universe looks the same at each point
(homogenous) and the Universe looks the same in all directions (isotropy). 

Because everything is moving away from everything else, we can possibly
retrace that in the distance past, everything originated from a state
where everything was much closer to each other. This is known as the
\textbf{Big Bang Cosmology.}
\end{defn*}

\subsection{Particles in the Universe}
\begin{defn*}
Total Energy of a particle:

\[
E_{total}^{2}=m^{2}c^{4}+p^{2}c^{2}
\]

where $m$ is the particle's rest mass, and $p$ the particle's momentum.
As mass is smaller, the particle moves more relativistically. 
\end{defn*}
Some particles in our Universe:
\begin{description}
\item [{Baryons}] Electrons, Protons and Neutrons. The Universe is charge
neutral so there is one electron for every proton. Also too heavy
to be moving at relativistic speeds
\item [{Radiation}] Photons $\gamma$, no mass with $E=hf$. Can knock
out an electron from an atom, or can scatter off a free electron (Thomson
scattering, or Compton Scattering when $hf\rightarrow m_{e}c^{2}$. 
\item [{Neutrinos}] extremely weakly interacting particles, produced, for
example, in radioactive decay. May posses non zero rest mass but treated
as massless, so always relativistic. Three types of neutrino, the
electron neutrino, muon neutrino and tau neutrino.
\item [{Dark~matter}] we don't know what this is\pagebreak{}
\end{description}
\vspace{10mm}

\section{Newtonian Gravity}

\subsection{Cheating, because I don't know GR}

Some equations are the same, and it's possible to cheat our way to
them without passing by GR. This is the case with the Friedmann equation. 

We start with two important equations and a theorem:

\noindent\fbox{\begin{minipage}[t]{1\columnwidth - 2\fboxsep - 2\fboxrule}%
Newtonian Equations:

\begin{equation}
\begin{aligned}F= & \frac{GMm}{r^{2}}\end{aligned}
\end{equation}

\begin{equation}
\begin{aligned}V= & -\frac{GMm}{r}\end{aligned}
\end{equation}

$F$ points in the direction that decreases $V$ because negative
potential energy means attraction. %
\end{minipage}}
\begin{thm}
From Classical Mechanics! Thanks Professor Cline! The force inside
a spherical shell is 0. The only relevant material is the mass inside
a concentric sphere through P.\label{Theorem1}
\end{thm}

The Friedmann equation describes the expansion of the Universe.

Consider an observer in a univerform expanding medium, with mass density
$\rho$ (mass per unit volume). We can consider any point to be it's
centre. A particle with mass $m$ with a distance $r$ compared to
the observer, will only feel a force from the mass within a radius
$r$ (\ref{Theorem1}). With the material having total mass $M=\frac{4\pi\rho r^{3}}{3}$,
it will contribute a force: 
\[
\begin{aligned}F= & \frac{GMm}{r^{2}}=\frac{G\left(\frac{4\pi\rho r^{3}}{3}\right)m}{r^{2}}=\frac{4\pi G\rho rm}{3}\end{aligned}
\]

With gravitational potential energy

\[
\begin{aligned}V= & -\frac{GMm}{r}=-\frac{G\left(\frac{4\pi\rho r^{3}}{3}\right)m}{r}=-\frac{4\pi G\rho r^{2}m}{3}\end{aligned}
\]

It's kinetic energy is given by 

\begin{align*}
T= & \frac{1}{2}m\dot{r}^{2}
\end{align*}

With the total energy $U$ given by:

\[
U=T+V
\]

with $U$ being a constant. Therefore

\begin{equation}
\begin{aligned}U= & \frac{1}{2}m\dot{r}^{2}-\frac{4\pi}{3}G\rho r^{2}m\end{aligned}
\label{eq:UFriedmann}
\end{equation}

This equation gives the evolution of the separation r between the
two particles.

We can change to comoving coordiantes, which describes comoving distarnce,
which factors out the expansion of the Universe, giving a distance
that does not change with time due to the expansion of space. The
real distance, corresponding to distance at a specific cosmological
time can change over time due to the expansion of the universe. 

\begin{equation}
\begin{aligned}\overrightarrow{r}= & a(t)\overrightarrow{x}\end{aligned}
\label{eq:expansionscale}
\end{equation}

$a(t)$ is the scale factor of the Universe. It measures the universal
expansion rate. It is a function of time alone due to the Cosmological
Principle.

Substituting equation (\ref{eq:UFriedmann}), with $\dot{x}=0$, as
objects are fixed in the comoving coordinates gives:

\[
\begin{aligned}U= & \frac{1}{2}m\dot{a}^{2}x^{2}-\frac{4\pi}{3}G\rho a^{2}x^{2}m\end{aligned}
\]
\marginpar{Remember that $\dot{r}=\dot{a}x+\dot{x}a=\dot{a}x$ from the chain
rule

so $\dot{r}^{2}=\dot{a}^{2}x^{2}$}

Multiplying each side by $\frac{2}{ma^{2}x^{2}}$ and rearranging
the terms then gives\marginpar{
\[
\protect\begin{aligned}U\times\left(\frac{2}{ma^{2}x^{2}}\right) & =\protect\\
\left(\frac{1}{2}m\dot{a}^{2}x^{2}-\frac{4\pi}{3}G\rho a^{2}x^{2}m\right)\times\frac{2}{ma^{2}x^{2}} & =\protect\\
\protect\\
\protect\end{aligned}
\]
 }:

\[
\begin{aligned}\left(\frac{\dot{a}}{a}\right)= & \frac{8\pi G}{3}\rho-\frac{kc^{2}}{a^{2}}\end{aligned}
\]

\end{document}
